\section{The wiring}
\label{sec:wiring}

\subsection{SPI between the boards}

Four wires plus power and ground, according to the contract's pin configuration. Two things are
worth spelling out:

\textbf{\code{VDDIO2} first.} \code{PORTC} is the MVIO domain and is fed from there. Without that
connection the port does nothing at all, and the symptom looks exactly like a broken SPI
implementation.

\textbf{A common ground is not optional.} Two boards without a common ground do not communicate,
and the symptom is incomprehensible.

\subsection{The CAN side}

The FPGA node's CAN side is a \textbf{single-wire open-drain bus} on three GPIO pins:
\code{tx_bus} (what the node drives), \code{bus_en} (whether it drives at all) and \code{rx_bus}
(the bus level coming in).

For a bus analyser to see the traffic, that single wire has to become a real
\textbf{differential} CAN bus. That is the transceiver's job, and its transmit input is to be
\textbf{low for a dominant bit}:

\begin{narrowtable}{@{}llll@{}}
\thead{\code{bus_en}} & \thead{\code{tx_bus}} & \thead{The node} & \thead{Transceiver input} \\ \midrule
0 & --- & does not drive & 1 (recessive) \\
1 & 0 & dominant & 0 \\
1 & 1 & recessive & 1 \\
\end{narrowtable}

How the two signals are combined onto the single wire is the hardware side's part of the wiring;
check it before you connect anything. A node whose \code{bus_en} has been ignored drives the bus
all the time and blocks everybody else.

\textbf{Termination: 120~\textohm{} at each end}, not in the middle and no more than two. A badly
terminated bus often works anyway on a short cable, and stops working precisely when you have
started to trust it.

\section{The ladder}
\label{sec:ladder}

Climb it in order. Every rung has a receipt of its own, and the point is to \textbf{never have
more than one unknown at a time}. Skip a rung and you are debugging four layers at once.

\begin{booktable}{@{}llL@{}}
\thead{Rung} & \thead{Do} & \thead{Receipt} \\ \midrule
0 & Power and ground & The MVIO status reports that \code{VDDIO2} is within range. \\
1 & \code{MOSI} to \code{MISO}, without the FPGA & You get back exactly the byte you sent. Try
  \code{0x00}, \code{0xFF}, \code{0xAA}, \code{0x55}. \\
2 & Connect the FPGA; write and read \code{TX_ID} & \code{0xFFFFFFFF} reads back as
  \code{0x000007FF}. \\
3 & Read \code{STATUS} repeatedly & \code{0x00000001} --- TX ready set, the others low. \\
4 & Connect the transceiver; run a \code{send()} & Activity on the bus; \code{STATUS} bit 0 goes
  low and comes back. \\
5 & Let a second node transmit & \code{receive()} returns \code{true} with the right content. \\
\end{booktable}

Two of the rungs deserve a comment.

\textbf{Rung 1 is \code{AvrSpiTransport}'s only test.} It is the only class in the project that no
automated test covers, since it talks to a physical register. If the loopback does not work, the
fault is in your code or in the configuration --- the FPGA is not connected and cannot be the
cause.

\textbf{Rung 2 is the decisive one.} The masking to eleven bits proves that you are talking to the
\emph{register bank} and not to an echo of your own wire. A loopback gives back what you sent; a
register bank gives back what it stored.

\section{Troubleshooting}
\label{sec:troubleshooting}

\begin{booktable}{@{}lLl@{}}
\thead{Symptom} & \thead{Likely cause} & \thead{Check} \\ \midrule
Nothing at all, no clock pulses & The pin mux is not set & Rung 1 \\
The port does nothing & \code{VDDIO2} not powered & The MVIO status \\
Every read gives \code{0x00000000} & \code{MISO} not connected, or the wrong pin & Swap
  \code{MOSI}/\code{MISO} and see whether it changes \\
Every read gives \code{0xFFFFFFFF} & \code{MISO} floating & The wiring \\
The reads are one byte out & The data register is not read in \code{transfer()} & \kapref{ch:avr} \\
Right for small numbers, not large & Sign extension in \code{readReg()} & \kapref{ch:byte} \\
The second transaction is ignored & Two transactions during one low \code{SS} & \kapref{ch:spi} \\
The frame arrives back to front & \code{TX_DATA_LO}/\code{HI} swapped & \kapref{ch:driver} \\
The same frame over and over & \code{RX_ACK} is not written & \kapref{ch:driver} \\
\code{hasError()} always \code{true} & \code{clearError()} writes \code{0x1} & \kapref{ch:driver} \\
The bus constantly dominant & A node's \code{bus_en} ignored & The wiring table \\
Works on a short cable, not a long & Termination & 120~\textohm{} at each end \\
\end{booktable}

\subsection{The procedure}

\begin{enumerate}
  \item Which rung was the last one to give the right receipt?
  \item What is the \textbf{one} thing that has been added since then?
  \item Formulate a hypothesis about what is wrong.
  \item Decide which measurement would settle the matter.
  \item Make the measurement.
\end{enumerate}

\begin{aside}
\note Swapping wires until it works is not troubleshooting, and it takes longer. The difference
between the two ways of working is about an afternoon per fault, and it is the most transferable
skill in the whole course.
\end{aside}

\subsection{Debugging somebody else's half}

A bug during bring-up can be in your code, in the other group's VHDL, or in a jumper wire. Three
rules make it manageable:

\begin{itemize}
  \item \textbf{Suspect your own half first.} Not out of humility, but because you can change it
    and measure on it.
  \item \textbf{Prove it before you report it.} "It does not work" is not a fault report. "Rung 2
    gives \code{0x000007FF} for \code{TX_ID} but \code{0x00000000} for \code{TX_DLC}" is.
  \item \textbf{Change one thing at a time, and write down what.} Two people debugging the same
    wiring at the same time make each other's measurements meaningless.
\end{itemize}
